% Packages, layout and theorem environments.
% Notation and shortcuts belong in macros.tex.

% --- Fonts and encoding ------------------------------------------------------
\usepackage[T1]{fontenc}
\usepackage{lmodern}
\usepackage[english]{babel}
\usepackage{microtype}
\usepackage{csquotes}
\usepackage{stix}

% --- Mathematics -------------------------------------------------------------
\usepackage{amsmath}
\usepackage{amssymb}
\usepackage{amsthm}
\usepackage{mathtools}
\usepackage{stmaryrd}
\SetSymbolFont{stmry}{bold}{U}{stmry}{m}{n} % stmaryrd has no bold series
\newcommand{\bmmax}{0} % bm reserves no bold math alphabets: TeX allows only 16 per version
\newcommand{\hmmax}{0}
\usepackage{bm}
\usepackage{mathpartir}
\usepackage{tikz-cd}
\usepackage{quiver}
\usepackage{eucal}

% --- Lists, colour, misc -----------------------------------------------------
\usepackage{enumitem}
\usepackage{xcolor}
\usepackage{booktabs}

% --- Bibliography ------------------------------------------------------------
\usepackage[backend=biber,style=alphabetic,maxbibnames=99,giveninits=true]{biblatex}

% --- Hyperlinks and cross-references (load last) -----------------------------
\usepackage[hidelinks]{hyperref}
\usepackage[capitalise,noabbrev]{cleveref}

% --- Theorem environments ----------------------------------------------------
% All numbered statements share one counter, numbered within the chapter.
% They are declared after cleveref so that \cref picks the right names.

\theoremstyle{plain}
\newtheorem{theorem}{Theorem}[chapter]
\newtheorem{proposition}[theorem]{Proposition}
\newtheorem{lemma}[theorem]{Lemma}
\newtheorem{corollary}[theorem]{Corollary}

\theoremstyle{definition}
\newtheorem{definition}[theorem]{Definition}
\newtheorem{example}[theorem]{Example}
\newtheorem{notation}[theorem]{Notation}

\theoremstyle{remark}
\newtheorem{remark}[theorem]{Remark}

\crefname{notation}{Notation}{Notations}

\numberwithin{equation}{chapter}
